

%% ===================================================================
\section{\vrev{종합 논의}}
\label{sec:ch5_discussion}
%% ===================================================================

제4장(연구 결과)의 교차분석의 분석 결과를 종합하여 DRTO 개별 시그모이드
바운딩 모델의 이론적 기여, 방법론적 함의, 그리고 한계를 논의한다.
교차분석에서 도출된 \jrev{\chg{주요 결과와 발견, 검증}은} 다음 세 가지 핵심 메시지로 통합된다.
\label{subsec:ch5_integration}

\chg{본 장의 논의는 \chref{ch:introduction}에서 제기한 다섯 가지 연구 목적에 대한 응답으로 구성된다.
서론은 관측성 단일 채널 모델을 관측성과 불확실성의 이중 채널 $\DRTO$ 모델로 확장하고,
정적 기준에서 관측성과 불확실성을 단계적으로 도입하는 네 개 운영 조건 체계(C0--C3)를 설계하며,
모델의 효과와 강건성, 회귀 설명력, 이중채널 비선형 효과, 전문가 타당성의 다섯 측면을 다층 검증하고,
관측성과 극단 불확실성이 결합할 때 나타나는 이중채널 비선형 효과를 규명하며, 6개 산업에 적용
가능한 실무 가이드라인을 제시하는 것을 목표로 제시하였다. 본 연구는 이들 목적이 제4장(연구 결과)로 어떻게
충족되는지를 5.2(학술적 기여)와 5.3(실무적 기여)로
나누어 논의하며, 제4장(연구 결과)의 모든 결과 기여와 활용 귀착은 5.1.1(제4장 연구 결과의 기여와 활용 종합 매핑)에 모두 매핑하였다. 또한 세 개의 핵심 메시지는 다음과 같다.}

\chg{첫째, 핵심 메시지는 관측성과 불확실성의 구조적 반대 효과이다.}
C0$\to$C1 전환(Cohen's $d = -1.180$)과 C1$\to$C2 전환($d = +1.871$)의
효과 크기는 모두 \jrev{\textbf{대효과(large effect)}} 수준이며 방향이 반대이다.
이는 개별 시그모이드 바운딩 모델에서 $E_{O,\text{bnd}} \leq 0$과
$E_{U,\text{bnd}} \geq 0$으로 설계된 구조적 반대 효과가
실험적으로 확인된 것이다.
\pdferr{2차 회귀}에서 C1은 $R^2 = 1.000$, C2는 $R^2 = 0.998$에 도달하여,
이 반대 효과가 \pdferr{2차 다항식}의 정밀도로 정량화됨을 보여준다.
실무적으로 이 발견은 BCM 전략 수립 시 관측성 투자(복구 시간 단축)와
불확실성 대비(안전 마진 확보) 사이의 trade-off를 명확히 인식해야
함을 시사한다.

\chg{둘째, 핵심 메시지는 평균의 비차별성과 극단의 극적 분기이다.}
C2--C3 비교에서 Cohen's $d = -0.238$(\jrev{소효과(small effect)})이지만,
$\mathrm{CVaR}_{99}$ 격차는 $+24.0\%$, Tail $k_O$ 계수 감쇠는 $0.52$배이다.
이 ``평균에서의 비차별성과 극단에서의 극적 분기'' 패턴은
불확실성의 분포 형태가 평균적 의사 결정이 아닌 극단 위험 관리에서
결정적으로 중요함을 정량적으로 입증한다.
P85.8 교차점은 \jrev{이 두 체계의 분기점(자연적 경계)}이며,
분할 회귀에서 Core($R^2 = 0.999$)와 Tail($R^2 = 0.710$)의
높은 설명력은 \jrev{이 분기의 통계적 정당성}을 확인한다.

\chg{셋째, 핵심 메시지는 전문가 검증에 의한 실무적 타당성 확인이다.}
5명의 전문가가 전체 평균 $4.46/5.0$으로 프레임워크를 평가하였으며
($t = 3.68$, $p = 0.011$), 15개 하위 주제 전체에서 $80\%$ 이상
동의가 확인되었다. 이 결과는 시뮬레이션에 기반한 이론적 발견이
BCM 실무자의 경험적 관찰과 부합함을 삼각 검증(triangulation)한 것이다.
특히 전문가들이 P85.8 교차점과 이중 채널 감쇠 메커니즘의 실무적
설명력을 높이 평가한 점은 이 발견의 이론적 신규성과 실무적
유용성을 동시에 확인한다.


\subsection{\chg{제4장 연구 결과의 기여와 활용 종합 매핑}}
\label{subsec:ch5_results_map}

\chg{본 항은 \chref{ch:results}에서 도출된 모든 결과가 학술적 기여와 실무적 활용 가운데 어디로 귀착되는지를
일람으로 제시한다. 이는 의미의 크고 작음과 무관하게 각 결과가 본 장의 어느 논의에서 성과로 활용되고
기여로 정리되는지를 명시함으로써, 결과와 논의 사이의 누락 없는 연결을 보장하기 위함이다. 학술적 기여는
\secref{sec:academic-contrib}에서, 실무적 활용은 \secref{sec:practical-contrib}에서 상세히 전개되며,
\tabref{tab:ch5-results-map}\refpo{tab:ch5-results-map}{은}{는} 그 매핑을 결과 영역별로 정리한 것이다.}

\begin{table}[htbp]
  \centering
  \caption{\chg{제4장 결과 영역별 학술적 기여와 실무적 활용 귀착}}
  \label{tab:ch5-results-map}
  \small
  \begin{tabularx}{\textwidth}{l X X}
    \toprule
    \textbf{\chg{결과 영역(4장 절)}} & \textbf{\chg{핵심 결과}} & \textbf{\chg{학술 기여 / 실무 활용 귀착}} \\
    \midrule
    \chg{관측성 효과 (4.1)} & \chg{$\bar\eta_{\text{C1}}=0.853$($-14.7\%$), $d=-1.180$, H1$\cdot$H3-a 100\%} & \chg{이론(5.2.1)$\cdot$실증(5.2.3) / 회귀식 산출(5.3.1)$\cdot$성숙도 진단(5.3.3)} \\
    \addlinespace
    \chg{가우시안 프리미엄 (4.2)} & \chg{$\bar\eta_{\text{C2}}=1.682$($+68.2\%$), $d=+1.871$, 프리미엄율 84.9\%} & \chg{이론(5.2.1)$\cdot$실증(5.2.3) / 회귀식 산출$\cdot$안전마진 해석(5.3.1)} \\
    \addlinespace
    \chg{파레토$\cdot$극단 위험 (4.3)} & \chg{$\bar\eta_{\text{C3}}=1.514$, P85.8 CDF 교차점, CVaR$_{99}$ $+24.0\%$, CVaR 에스컬레이션} & \chg{이론(5.2.1) / \textbf{이중채널 자원배분$\cdot$극단 대비(5.3.2)}} \\
    \addlinespace
    \chg{회귀$\cdot$간명모형 (4.4)} & \chg{$R^2$ 진행(C1 1.000/C2 0.998/C3 0.858), 분할 Core 0.999, 4변수 간명모형, 지배변수 전환} & \chg{방법론(5.2.2) / \textbf{회귀식 DRTO 산출 도구(5.3.1)}$\cdot$산업별 투자우선순위(5.3.4)} \\
    \addlinespace
    \chg{이중채널 감쇠 (4.10)} & \chg{Tail $k_O$ $-0.415$ vs Core $-0.798$($0.52$배), Core 선형/Tail 포화 메커니즘} & \chg{이론(5.2.1) / \textbf{이중채널 스트레스테스트$\cdot$자원배분(5.3.2)}} \\
    \addlinespace
    \chg{\m{효과 크기}$\cdot$분포구조 (4.7--4.9)} & \chg{경로별 Cohen's $d$, 순효과 격차(가우시안 1.110 vs 파레토 0.650), 백분위$\cdot$PDF 교차, 회귀계수 직접감쇠($k_O$ 0.75배$\cdot$U 0.45배)} & \chg{이론(5.2.1)$\cdot$실증(5.2.3) / 산업별 환경계수$\cdot$투자효과 비교(5.3.4)} \\
    \addlinespace
    \chg{강건성$\cdot$기반검증 (4.5--4.6)} & \chg{기술통계 질적전환, $n{=}10{,}000$ 적정성, $10^8$ 다중시드(H6$\cdot$H7), Golden Compare 47점(L0)} & \chg{방법론(5.2.2)$\cdot$실증(5.2.3) / 재현성 신뢰 근거(5.3.3)} \\
    \addlinespace
    \chg{전문가 검증 (4.11)} & \chg{$4.46/5.0$($p{=}0.011$, $d{=}1.65$), CVI $0.985$, $\alpha{=}0.89$, 인터뷰 15/15 $\geq80\%$(H9$\cdot$H10)} & \chg{실증(5.2.3) / 실무 수용성$\cdot$현장 통찰(5.3 전반)} \\
    \addlinespace
    \chg{교차 정합성 (4.12)} & \chg{\citet{lee2026drto} 매개변수 승계 일관성, 303개 교차검증 PASS, 한계-해결 연쇄} & \chg{방법론(5.2.2)$\cdot$실증(5.2.3) / 연구 신뢰성 보증} \\
    \addlinespace
    \chg{가설 종합 (4.13)} & \chg{H1--H10 전원 채택, \frev{분포 환경 적합성} 종합 명제, $\kappa{=}1.2$ 잠정 채택} & \chg{실증(5.2.3)$\cdot$한계(5.4) / $\kappa$ 기본값 적용 근거(5.3)} \\
    \bottomrule
  \end{tabularx}
\end{table}

\chg{이 매핑이 보이는 핵심은, 제4장(연구 결과)의 결과가 학술적으로는 이론, 방법론, 실증의 세 기여로, 실무적으로는
회귀식 기반 산출, 이중채널 기반 자원배분, 성숙도와 에스컬레이션 운영, 산업별 적용의 네 도구로 빠짐없이
이어진다는 점이다. 특히 본 연구의 두 핵심 발견인 \textbf{회귀식(간명모형)}과 \textbf{이중채널 감쇠
메커니즘}은 각각 5.3(실무적 기여)의 독립된 실무 도구로 구체화되며, 이는 이론적 발견이
실무 의사결정으로 내려오는 경로를 명시한다.}


